% ============================================================
\section{Literature Review}
\label{sec:literature}
% ============================================================

The strategy builds on five interconnected bodies of evidence. We set them out in the order they inform the design.

The core signal is momentum. \citet{jegadeesh1993} document that past equity winners outperform past losers over intermediate horizons, and \citet{asness2013} extend the evidence across asset classes and geographies. \citet{moskowitz2012} establish the time-series counterpart: an asset's own past return predicts its future return across equities, bonds, currencies, and commodities. A survey by \citet{jegadeesh2023} confirms both effects remain robust three decades on. A practical difficulty in applying momentum is that no single lookback horizon dominates consistently across assets and periods: short windows capture reversals rather than continuation, while very long windows can capture mean-reversion. The 13612U signal in Rule~1 addresses this by averaging 1-, 3-, 6-, and 12-month excess returns, blending horizons known to carry cross-sectional and time-series momentum in approximately equal measure \citep{moskowitz2012, keller2017}. The pre-filter in Rule~2 restricts candidates to the top-$7$ assets by this score before the absolute threshold is applied: concentrating the eligible set on the highest-momentum names reduces exposure to assets near the boundary of positive momentum and has been shown to improve return and drawdown outcomes in multi-asset rotation \citep{faber2007, keller2017}. Several papers have translated the momentum insight into ETF-based rotation strategies: \citet{faber2007} shows that a moving-average rule materially reduces drawdowns in a diversified asset portfolio; \citet{blitz2008} apply value and momentum across asset classes; \citet{clare2016} combine trend following and risk-parity ideas in global allocation; and \citet{bellu2020} propose a momentum framework with explicit cash protection in downturns, the closest published precedent to our strategy. More recent work by \citet{wouassom2022} and \citet{vukovic2023} documents momentum in international equity markets. The present paper shares this lineage but differs in focus: rather than evaluating a strategy as a whole, it attributes out-of-sample performance rule by rule.

\citet{antonacci2014} formalises the combination of cross-sectional and absolute momentum as \emph{dual momentum}: assets are first ranked relative to peers, then screened against a positive-trend threshold. The absolute filter is the cash-protection mechanism, preventing the portfolio from holding assets in a downtrend regardless of relative rank. \citet{antonacci2012} and \citet{antonacci2015} provide supporting empirical evidence. Our Rule~3 retains only assets whose 13612U score exceeds the cash rate ($\theta = 2\%$), directly implementing this filter, and the breadth-based defensive switch extends the same logic to the portfolio level: when momentum collapses broadly, the portfolio exits risky assets entirely.

\citet{keller2017} extend the absolute-momentum concept to a \emph{breadth} signal: the fraction of assets with positive momentum serves as a regime indicator for the Vigilant and Hybrid Asset Allocation frameworks. Breadth-based switching is well motivated by \citet{daniel2016}, who show that momentum crashes cluster in panic states where most assets deteriorate simultaneously, and by \citet{barroso2015}, who document that momentum is most vulnerable precisely when volatility is elevated. Triggering the defensive switch through breadth collapse uses the same signals that drive offensive allocation, making regime detection endogenous and fully auditable.

When selecting a small portfolio from a large eligible set, the covariance structure matters. \citet{markowitz1952} establishes diversification as the mechanism for reducing portfolio risk; in practice, however, estimating the full covariance matrix for a large universe introduces substantial estimation error. \citet{demiguel2009} demonstrate that naive equal-weight ($1/N$) allocation is hard to beat net of this error, a result that motivates our equal-weight design. Once an absolute momentum filter has been applied, the eligible assets share similar expected returns (as we verify on the training data), so the only remaining lever for portfolio construction is their correlation structure: selecting the lowest-correlation triplet minimises portfolio variance while bypassing matrix inversion. This rule-based, estimation-light approach is the portfolio-theoretic foundation of Rule~4. Neither \citet{faber2007} nor \citet{keller2017} incorporate a correlation-based selection criterion within the eligible set; once assets pass the absolute filter, both frameworks weight them equally without regard to their pairwise correlations. The minimum-correlation triplet is the feature that most directly differentiates the present strategy from this prior work, and characterising its precise role through a skill-versus-luck attribution is the paper's central question.

The flexibility of rule-based backtests makes them vulnerable to data snooping. \citet{sullivan1999} and \citet{white2000} formalise this concern for technical trading rules; \citet{harvey2016} document pervasive multiple-testing bias in the cross-section of expected returns; and \citet{bailey2014} propose the Deflated Sharpe Ratio as a multiple-trial correction. \citet{yang2019} apply these cautions directly to tactical allocation rules. The existing literature typically delivers a snooping-corrected verdict on a strategy as a whole. It less often asks which rule inside the strategy carries the edge. We bring the two together, attributing performance rule by rule through a skill-versus-luck permutation test and a nested spanning-alpha ladder. The combination localises the edge and, at the same time, offers a constructive defence against the snooping critique.
